% Figure: induction vs strong induction --- what each step is allowed to lean on.
% Source for Figure~\ref{fig:reachback}. TikZ only; regenerable by compiling main.tex.
\begin{tikzpicture}[
  dot/.style={circle,draw,inner sep=1.4pt,minimum size=5.2mm,font=\footnotesize},
  bd/.style={circle,draw=accent,fill=accentsoft,inner sep=1.4pt,minimum size=5.2mm,font=\footnotesize},
  arr/.style={->,accent},
  sarr/.style={->,accent,bend left=24},
  lab/.style={font=\footnotesize\sffamily,text=slate},
]
% ---- ordinary induction ----
\node[lab,anchor=east] at (-0.3,0) {induction};
\foreach \i in {1,...,7}
  \node[dot] (a\i) at (\i*0.72,0) {\i};
\node[bd] at (0.72,0) {1};
\foreach \i [evaluate=\i as \j using int(\i+1)] in {1,...,6}
  \draw[arr] (a\i.east) -- (a\j.west);
\node[lab,anchor=west] at (5.8,0) {each step leans on the previous case only};
% ---- strong induction ----
\node[lab,anchor=east] at (-0.3,-2.1) {strong};
\foreach \i in {1,...,7}
  \node[dot] (b\i) at (\i*0.72,-2.1) {\i};
\foreach \i in {1,2,3}
  \node[bd] at (\i*0.72,-2.1) {\i};
% arrows from all earlier cases into n+1 (show into 4 and 7)
\foreach \i in {1,2,3}
  \draw[sarr] (b\i.north) to (b4.north);
\foreach \i in {4,5,6}
  \draw[sarr] (b\i.north) to (b7.north);
\node[lab,anchor=west] at (5.8,-2.1) {each step may lean on \emph{any} earlier case};
\node[lab,anchor=north] at (1.44,-2.5) {base window};
\draw[rulegray] (0.44,-2.5) -- (2.44,-2.5);
\end{tikzpicture}
